\chapter{Two traditions}
\label{chap:background}

The model of this thesis sits at the meeting of two bodies of work that have, for the most
part, developed apart. One is the econophysics of firm growth and of the size distributions of
firms and wealth: it has the empirical facts and the compositional models built to explain them,
and a separate line of scale-invariant growth models that produce heavy-tailed sizes. The other
is the statistical physics of large disordered dynamical systems, the random Lotka--Volterra and
replicator equations of theoretical ecology, which has the machinery of a persistently
fluctuating many-body state and its mean-field theory but has not been turned on firm growth.
Neither tradition on its own yields the result of the following chapters. This chapter sets out
what each provides, so that the model, and what in it is new, can be read against them.

\section{Compositional and granular firm-growth models}
\label{sec:bg-granular}

The granular argument of Chapter~\ref{chap:intro} says that a firm's growth statistics
come from its internal composition: the firm is an aggregate of many smaller units, each
fluctuating on its own. \citet{stanley1996} and \citet{amaral1997} first documented the
anomalous size--variance scaling. \citet{sutton2002} derived it by treating a firm as a
random partition into sub-units, with all partitions equally likely.
\citet{wyartbouchaud2003} obtained it from sub-units with power-law distributed sizes.
\citet{gabaix2011} applies the same logic one level up: when the firm-size distribution is
heavy-tailed, the largest firms do not average out, and their idiosyncratic shocks survive
as aggregate fluctuations. \citet{bottazzi2006} give the distributional companion, a
tent-shaped growth distribution generated by cumulative allocation of business opportunities. What these models share,
and what \citet{moran2024} call their ``island'' character, is that firms do not interact: each
firm's statistics come from its own internal composition, and the correlations between firms
that a network would carry are absent by construction. \citet{moran2024} then show that the
data break the further predictions of this picture (the thin tail of the rescaled volatility
distribution, the persistence of fat tails at all sizes, and the order-dependence of the
higher volatility moments) and conjecture a correlation that
grows with firm size, of the kind supply chains and shared customers would produce. A recent
working paper by Fontanelli, Napoletano and Secchi \citep{fontanelli2025} reopens even the first
fact, reporting a
size--volatility relation that steepens for small firms and flattens for large ones rather than
holding a single exponent. The compositional tradition thus arrives at the
need for a between-firm mechanism it does not itself contain.

\section{Scale-invariant growth and mean-coupled dynamics}
\label{sec:bg-scaleinv}

A separate econophysics line builds the scale invariance directly into the dynamics.
\citet{bouchaudmezard2000} model wealth by
\begin{equation}
    \dot W_i = \eta_i(t)\,W_i + \sum_{j} J_{ij} W_j - \sum_{j} J_{ji} W_i,
    \label{eq:bg-bm}
\end{equation}
built, in their words, so that ``since the unit of money is arbitrary'' the dynamics are
unchanged when all $W_i$ are multiplied by a common factor. In the mean-field case
$J_{ij} = J/N$ this reduces to $\dot W_i = \eta_i W_i + J(\bar W - W_i)$, each agent exchanging
with the population average $\bar W$, and the aggregate grows exponentially, $\bar W(t) \propto
e^{(m+\sigma^2)t}$. The Lotka--Volterra models of wealth and firm market values of the
Solomon school \citep{malcai2002} likewise couple each unit to the mean, $w_i(t+1) =
[1+\lambda(t)]w_i(t) + a\,\bar w(t) - c\,w_i(t)\bar w(t)$, though there the coupling is
quadratic. Both models produce stationary Pareto or power-law size
distributions, consistent with the empirical heavy tail documented by \citet{axtell2001}, and in both the fluctuations
are supplied by an external multiplicative noise $\eta_i$ rather than generated by the
interactions. The present model inherits its scale invariance and the mean-coupled form of its
immigration term from this line, and makes the
self-limitation and the disordered interactions relative as well, so that the growth is
endogenous and a persistently fluctuating state, rather than a static size distribution, carries
the firm-growth statistics.

\section{Firm volatility on production networks}
\label{sec:bg-network}

In production-network models of the macroeconomy, each firm is linked to its suppliers
and customers, and a shock to one firm travels along these links. Linked firms fluctuate
together, so their shocks do not cancel in the aggregate, and the fluctuations of the
whole economy depend on the structure of the network, which is exactly the kind of
between-firm structure the granular models leave out
\citep{acemoglu2012, carvalho2019}. Brought down to the single firm, \citet{herskovic2020}
build a customer--supplier network model in which a firm's volatility is set by its position in
the network and the concentration of its customer base, and show that the dispersion of the
firm-size distribution is the common factor in firm volatilities. This is the closest existing
account to the mechanism of the present thesis, but it is a static,
equilibrium model of the \emph{level} of firm volatility, whereas at issue here is the
\emph{size-conditioned} structure (the exponent of the size--volatility law and the
order-dependence of its higher moments) and its origin in a fluctuating dynamics, which no
account in this line addresses.

\section{Disordered Lotka--Volterra and its fluctuating phase}
\label{sec:bg-glv}

The dynamics the thesis uses come from a third body of work, the statistical physics of large
disordered interacting systems. May's random-interaction argument provides the classical
complexity--stability benchmark \citep{may1972}. The random replicator equation \citep{diederich1989,
opperdiederich1992} and the disordered generalized Lotka--Volterra model of theoretical ecology
\citep{bunin2017} describe $N$ units with random pairwise couplings, and a control parameter,
the disorder strength $\sigma$, sets their collective behaviour. Below a threshold the dynamics
relax to a fixed point; above it, related models can exhibit persistently fluctuating states. In a
spatially heterogeneous metacommunity these states can be chaotic \citep{roy2020}, and a single
disordered Lotka--Volterra community with weak immigration displays chaotic turnover
\citep{mallmin2024}. That fluctuating
phase is the property the compositional and mean-coupled models above lack and the thesis
needs: a state in which every unit keeps fluctuating without external noise, the fluctuations
generated by the quenched disorder of the couplings alone \citep{roy2020}.

Related analyses use a dynamical mean-field theory (DMFT) that, for the models and assumptions
they study, reduces the $N$-body dynamics to an effective single-unit process in a self-consistent
field \citep{galla2018, roy2020}; Appendix~\ref{app:dmft} carries out the corresponding
derivation for the relative model. On a
heterogeneous graph the reduction becomes degree-resolved: \citet{aguirrelopez} develops a
heterogeneous DMFT in which the field a unit feels carries an amplitude set by its degree,
$\sqrt{k_i/C}$, though the analysis is of the fixed points of the non-growing model rather than
of the fluctuating phase. In the following chapters, the graph's degree distribution enters the firm-growth
statistics through this field. Two
features the thesis relies on have close counterparts in recent work on the fluctuating phase:
the confinement of abundant units by their own self-limitation and the intermittent turnover of
which unit is large \citep{mallmin2024}, and the single-unit process with a migration floor, a
finite-correlation-time field, and a self-limitation whose strength grows with abundance
\citep{arnouxbunin2024}. This literature has not carried the fluctuating phase onto a
scale-free graph, made the interactions relative so that the aggregate grows, or read off the
size-conditioned firm-growth statistics. That is the work of the next two chapters,
where the degree heterogeneity turns out to generate the multiscaling the compositional
tradition cannot.
